\subsection{Operators with compact resolvent}

In this number, E denotes a complex Hilbert space of infinite dimension.

PROPOSITION 19. --- Let u be a self-adjoint partial operator on E. The following conditions are equivalent:

\begin{enumerate}
\def\labelenumi{(\roman{enumi})}
\item
  There exists an orthonormal basis \(\mathrm{B} = (e_{j})_{j \in \mathrm{J}}\) of E such that u is diagonal in the basis B and the absolute value of the family of eigenvalues of u tends to infinity along the filter of complements of finite subsets of J;
\item
  For every \(\lambda\) belonging to the resolvent set of \(u\) , the resolvent \(\mathrm{R}(u,\lambda)\) is compact;
\item
  There exists a complex number \(\lambda\) belonging to the resolvent set of \(u\) such that the resolvent \(\mathrm{R}(u,\lambda)\) is compact;
\item
  The spectrum of u coincides with the sensitive spectrum of u.
\end{enumerate}

Suppose that (i) is satisfied and let \((\lambda_{j})_{j\in\mathrm{J}}\) be the family of eigenvalues of u. Let \(\mu\) be the counting measure on J. The spectrum of u is the image \(\mu\) -essential of the family \((\lambda_{j})\) (prop. 22 of IV, p. 252); it is the set of values of this family. For every \(\lambda\notin\mathrm{Sp}(u)\) , the resolvent \(\mathrm{R}(u,\lambda)\) is diagonal in the basis B and the family of its eigenvalues is \(((\lambda-\lambda_{j})^{-1})_{j\in\mathrm{J}}\) (loc. cit.). Since this family converges to 0, the endomorphism \(\mathrm{R}(u,\lambda)\) is compact (prop. 2, (iii) of IV, p. 148). This proves that (i) implies (ii).

Since the resolvent set of \(u\) is nonempty (cf. prop. 17 of IV, p. 248), properties (ii) and (iii) are equivalent by formula (8) of IV, p. 245 and proposition 3 of III, p. 5.

Since \(\mathrm{R}(u,\lambda)\) is normal for \(\lambda\notin\mathrm{Sp}(u)\) (prop. 16 of IV, p. 247), condition (iii) implies (iv) by prop. 15 of IV, p. 247 and prop. 5 of III, p. 90.

Finally, let us prove that (iv) implies (i). Let \(\mathcal{O}\) be the set of orthonormal subsets of E consisting of eigenvectors for \(u\) . The set \(\mathcal{O}\) , ordered by inclusion, is of finite character (E, III, p. 34, def. 2) since O belongs to \(\mathcal{O}\) if and only if the sets formed of at most two elements of O belong to \(\mathcal{O}\) . By E, III, p. 35, th. 1, there exists a maximal element O of \(\mathcal{O}\) . Let F be the closed subspace of E generated by O. For \(e \in O\) , there exists a unique \(\lambda(e) \in \mathbf{R}\) such that \(u(e) = \lambda(e)e\) (prop. 17 of IV, p. 248).

The set of values of the map \(\lambda\) from O to \(\mathbf{R}\) coincides with the spectrum of \(u\) . Indeed, on the one hand this set is contained in the spectrum of \(u\) and on the other hand, if there exists \(\lambda_0 \in \mathrm{Sp}(u)\) which is not a value of \(\lambda\) , then \(\lambda_0\) is an eigenvalue of \(u\) by hypothesis. There exists a vector \(e \in \mathrm{dom}(u)\) of norm 1 such that \(u(e) = \lambda_0 e\) and \(e \notin \mathrm{O}\) ; the subset \(\mathrm{O} \cup \{e\}\) is orthonormal (since eigenspaces of \(u\) corresponding to distinct eigenvalues are orthogonal by assertion \(b\) ) of the cor. of prop. 17 of IV, p. 248); it belongs to \(\mathcal{O}\) , contradicting the maximality of O.

Suppose that \(F \neq E\) . Then \(F^{\circ}\) is nonzero. The endomorphism \(\mathrm{R}(u,i)\) of E is normal (prop. 16 of IV, p. 247). It leaves the subspace \(F^{\circ}\) of E stable (lemma 4 of I, p. 135). Let v be the endomorphism of \(F^{\circ}\) deduced from \(\mathrm{R}(u,i)\) by passing to subspaces. It is a continuous and normal endomorphism of \(F^{\circ}\) (loc. cit.) whose spectrum is contained in that of \(\mathrm{R}(u,i)\) . Since \(F^{\circ}\) is nonzero, the spectrum of v is nonempty (cor. 1 of I, p. 26). Moreover, the spectrum of v is not reduced to 0, since the normal endomorphism v is nonzero (example 1 of I, p. 110). Let \(s \in \mathrm{Sp}(v) - \{0\}\) . Since s belongs to the spectrum of \(\mathrm{R}(u,i)\) , there exists \(e \in O\) such that \(s = (i - \lambda(e))^{-1}\) , and s is an eigenvalue of \(\mathrm{R}(u,i)\) (prop. 15, a) of IV, p. 247). Since s is nonzero, it is an isolated point of \(\mathrm{Sp}(\mathrm{R}(u,i))\) by hypothesis, hence also of \(\mathrm{Sp}(v)\) . Thus, s is an eigenvalue of v (prop. 5, c) of I, p. 134). Let \(e \in F^{\circ}\) be an eigenvector of norm 1 of v; it is also an eigenvector of u (prop. 15, b) of IV, p. 247), and the set \(O \cup \{e\}\) contradicts the fact that O is maximal in O. Hence F = E.

The family of elements of O is therefore an orthonormal basis of E consisting of eigenvectors of u. The spectrum of u is closed in R, and the sensitive spectrum is discrete; since these sets coincide, the set of elements \(\lambda \in \mathrm{Sp}(u)\) such that \(|\lambda| \leqslant R\) is compact, hence finite, for every R > 0. Thus the family of absolute values of the eigenvalues of u tends to infinity along the filter of complements of finite subsets of O. Hence (iv) implies (i).

DEFINITION 12. --- Let u be a self-adjoint partial operator on E. One says that u has compact resolvent if the equivalent conditions of Prop. 19 are satisfied.

PROPOSITION 20. --- Let u be a self-adjoint partial operator on E. Then u has compact resolvent if and only if the canonical injection j of the Hilbert space \(E_{u}\) into E is compact.

Suppose that u has compact resolvent. There exists \(\lambda \notin \operatorname{Sp}(u)\) such that \(\mathrm{R}(u, \lambda)\) is compact (prop. 19, (iii)). Let B be the unit ball of the Hilbert space \(E_{u}\) . Since u is a continuous linear map from \(E_{u}\) into E, the subset \(\mathrm{B}' = (\lambda 1_{\mathrm{E}} - u)(\mathrm{B})\) of E is bounded. Since \(\mathrm{R}(u, \lambda)\) is compact, the subset \(\mathrm{B} = \mathrm{R}(u, \lambda)(\mathrm{B}')\) is relatively compact in E (Remark 1 of III, p. 2). This proves that j is compact.

Conversely, suppose that the linear map \(j\) is compact. The complex number \(i\) belongs to the resolvent set of \(u\) (prop. 17 of IV, p. 248). We have \(u \circ \mathrm{R}(u, i) = -1_{\mathrm{E}} + i\mathrm{R}(u, i)\) . Let B be the unit ball of E. For every \(x \in \mathrm{B}\) , we have

\[
\| u \circ \mathrm{R} (u, i) (x) \| = \| - x + i \mathrm{R} (u, i) (x) \| \leqslant 1 + \| \mathrm{R} (u, i) \|,
\]

and consequently

\[
\begin{array}{r l} & {\| \mathrm{R} (u, i) x \| _ {u} ^ {2} = \| \mathrm{R} (u, i) x \| ^ {2} + \| u \circ \mathrm{R} (u, i) x \| ^ {2}} \\ & {\qquad \leqslant \| \mathrm{R} (u, i) \| ^ {2} + (1 + \| \mathrm{R} (u, i) \|) ^ {2}.} \end{array}
\]

The image C of B under \(R(u,i)\) is therefore bounded in \(E_{u}\) . Since j is compact by hypothesis, the set \(C = j(C)\) is relatively compact in E (III, loc. cit.). Consequently, the resolvent \(R(u,i)\) is compact and u has compact resolvent (prop. 19, (iii)).

COROLLARY. --- Let q be a closed positive partial form on E and u the positive self-adjoint operator representing q. The following conditions are equivalent:

\begin{enumerate}
\def\labelenumi{(\roman{enumi})}
\item
  The partial operator u has compact resolvent;
\item
  The positive endomorphism \(\sqrt{(1_{\mathrm{E}} + u)^{-1}} = (1_{\mathrm{E}} + u)^{-1 / 2}\) of E is compact;
\item
  The canonical injection j of the Hilbertian space \(E_{q}\) (def. 9 of IV, p. 292) in E is compact.
\end{enumerate}

Since u is positive, the real number -1 belongs to the resolvent set of u (prop. 17 of IV, p. 248), hence the positive endomorphism \(v = \sqrt{(1_{\mathrm{E}} + u)^{-1}}\) is defined, and we have \(v = (1_{\mathrm{E}} + u)^{-1/2}\) by the functional calculus.

The endomorphism \(v\) is compact if and only if \(v^2 = (1_{\mathrm{E}} + u)^{-1}\) is compact (prop. 6 of III, p. 91), that is, if and only if \(u\) has compact resolvent (prop. 19, (iii)). This proves that conditions (i) and (ii) are equivalent.

Suppose that the endomorphism \(v\) is compact. Let B be the unit ball of the Hilbert space \(\mathrm{E}_q\) . Since \((1_{\mathrm{E}} + u)^{1/2}\) is a continuous linear map from \(\mathrm{E}_q\) into E (cor. of th. 3 of IV, p. 294), the subset \(\mathrm{B}' = (1_{\mathrm{E}} + u)^{1/2}(\mathrm{B})\) of E is bounded, hence the subset \(j(\mathrm{B}) = \mathrm{B} = v(\mathrm{B}')\) is relatively compact in E (Remark 1 of III, p. 2). This proves that \(j\) is compact. Thus (ii) implies (iii).

The linear map \(\widetilde{v}\colon x\mapsto(1_{\mathrm{E}}+u)^{-1/2}(x)\) from E into \(\mathrm{E}_{q}\) is well defined and isometric (corollary of Th. 3 of IV, p. 294). Since \(v=j\circ\widetilde{v}\) , condition (iii) implies (ii) (Prop. 3 of III, p. 5).

Example. --- Let \(n \in N\) . Let U be an open subset of \(R^{n}\) . We equip U with Lebesgue measure, denoted \(\mu\) . Let \(\Delta\) be the scalar differential operator \(-\sum_{i=1}^{n}\partial_{i}^{2}\) on U. The partial operator \(\Delta_{-}\) with domain \(\mathcal{D}(U)\) defined by \(\varphi \mapsto \Delta(\varphi)\) is closable (prop. 13 of IV, p. 242) and symmetric (IV, p. 243). It is positive, since for every \(\varphi \in \mathcal{D}(U)\) , we have

\[
\begin{array}{r l} & {\langle \varphi | \Delta_ {-} (\varphi) \rangle = \int_ {\mathrm{U}} \overline {{\varphi}} \Delta (\varphi) d \mu = - \sum_ {i = 1} ^ {n} \int_ {\mathrm{U}} \overline {{\varphi}} \partial_ {i} ^ {2} \varphi d \mu} \\ & {\qquad = \sum_ {i = 1} ^ {n} \int_ {\mathrm{U}} \overline {{\partial_ {i} \varphi}} \partial_ {i} \varphi d \mu = \int_ {\mathrm{U}} \sum_ {i = 1} ^ {n} | \partial_ {i} \varphi | ^ {2} d \mu \geqslant 0.} \end{array}
\]

We denote by \(\Delta_{D}\) the Friedrichs extension of the positive symmetric partial operator \(\Delta_{-}\) (IV, p. 295, example); it is a Laplacian on U, called the Dirichlet Laplacian on U.

Let q be the positive partial form associated with \(\Delta_{D}\) . The domain of q is the Hilbertian space completion of \(\mathcal{D}(U)\) for the positive-definite Hermitian form

defined by

\[
(\varphi_ {1}, \varphi_ {2}) \mapsto \int_ {\mathrm{U}} \overline {{\varphi}} _ {1} \varphi_ {2} + \sum_ {i = 1} ^ {n} \int_ {\mathrm{U}} \overline {{\partial_ {i} \varphi_ {1}}} \partial_ {i} \varphi_ {2}
\]

for all \((\varphi_{1},\varphi_{2})\in\mathcal{D}(\mathrm{U})\times\mathcal{D}(\mathrm{U})\) . In other words, the domain of \(q\) is the Sobolev space \(\mathrm{H}_{0}^{1}(\mathrm{U})\) (No. 14 of IV, p. 221).

*Suppose that U is bounded. The canonical injection of \(\mathrm{H}_0^1 (\mathrm{U})\) into \(\mathrm{L}^2 (\mathrm{U})\) is compact; the Dirichlet Laplacian on U is therefore an operator with compact resolvent (corollary of Prop. 20). Since the Hilbert space \(\mathrm{H}_0^1 (\mathrm{U})\) is of countable type (prop. 20 of IV, p. 222), and since the image of \(\Delta_{\mathrm{D}}\) is of infinite dimension, there exists an increasing sequence \((\lambda_n)_{n\geqslant 0}\) of real numbers tending to \(+\infty\) and an orthonormal basis \((f_n)_{n\in \mathbf{N}}\) of \(\mathrm{L}^2 (\mathrm{U})\) whose elements belong to the domain of \(\Delta_{\mathrm{D}}\) , such that \(\Delta_{\mathrm{D}}(f_n) = \lambda_n f_n\) for all \(n\in \mathbf{N}\) . One can prove (“Weyl’s law”) that when T tends to \(+\infty\) , we have

\[
\sum_{\substack{n\in \mathbf{N}\\ \lambda_{n}\leqslant \mathrm{T}}}1\sim \frac{c_{n}}{(2\pi)^{n}} m\mathrm{T}^{n / 2}
\]

where \(c_{n} = \pi^{n/2}/\Gamma(1 + n/2)\) is the volume of the unit ball in \(R^{n}\) (INT, V, p. 101, § 8, no. 7) and \(m > 0\) is the Lebesgue measure of U.*

